\documentclass[portuguese,brazilian,12pt,final]{article}
\usepackage[brazil]{babel}      % Pacotes da lingua portuguesa
\usepackage[utf8]{inputenc}     % com os caracteres com ç e acentuação
                                % e hifenização em português


\usepackage{isomath}       % Pacote matemático
\usepackage{natbib}        % habilita novos comandos de citação
\usepackage{amsmath}       % Pacote matemático
\usepackage{graphicx}      % Para inserir figuras
\usepackage{indentfirst}   % Para identar o primeiro paragrafo da seção



% Definição de novos comandos:
\newcommand{\Dx}{\Delta x}
\newcommand{\Dy}{\Delta y}
\newcommand{\imax}{i_{\max}}
\newcommand{\jmax}{j_{\min}}





% Comandos de formatação da página
\setlength{\baselineskip}{16.0pt}
\setlength{\parskip}{3pt plus 2pt}
\setlength{\parindent}{20pt}
\setlength{\oddsidemargin}{0.5cm}
\setlength{\evensidemargin}{0.5cm}
\setlength{\marginparsep}{0.75cm}
\setlength{\marginparwidth}{2.5cm}
\setlength{\marginparpush}{1.0cm}
\setlength{\textwidth}{150mm}



% Acima é o preâmbulo

%------------------------
% aqui começa o documento
%------------------------
\begin{document}

\begin{center}
{\large \LaTeX para Iniciantes\\ Exemplo} \\ % \\ = nova linha
Última compilação em \today
\end{center}

\section{Introdução}

\LaTeX\ parece mais difícil do que realmente é. É quase tão fácil como $\pi$.

Veja só como é fácil escrever símbolos especiais como: $\alpha$,
$\beta$, $\gamma$,
$\delta$, $\sin x$, $\hbar$, $\lambda$,~$\ldots$
É possível também fazer subescritos $A_{x}$, $A_{xy}$
e sobrescritos  $e^x$, $e^{x^2}$ e
$e^{a^b}$.
Nós usaremos \LaTeX, que é baseado no \TeX\ porém com mais opções de macros de alto nível.

Os pacotes mais comuns e recomendados como utilização padrão em todos os documentos são:

\emph{isomath} e \emph{amsmath}: São pacotes matemáticos.

\emph{natbib}: Pacote que expande os comandos e citação

\emph{graphicx}: Pacote para inserir figuras

\emph{color}: Pacote de cores



Para inserir um novo parágrafo, deixe uma ou mias linhas em branco.
Note que todas as formulas estão entre
\$ que signifíca \textit{modo matemático inline}.

A fonte padrão é Computer Modern, que inclui as variações \textit{itálico},
\textbf{negrito},
\textsl{itálico2}, \textsc{letras maiúsculas pequenas}, \textsf{sem serife}, \textrm{estilo romano} e \texttt{espaçamento único}.

É possível também alterar o tamanho das letras, como por exemplo: \tiny esse texto é minúsculo, \small esse texto é pequeno, \Large esse texto é Grande, \Huge esse texto é Gigantesco, \normalsize tamanho normal.

O controle de espaços no texto é muito simples, como por exemplo a partir desse ponto. \hspace{2cm} existe um espaço de 2cm. e abaixo, tem mais 1cm de espaço entre essa e a proxima linha.

\vspace{1cm}

Proxima linha...

\section{Listas}

Alguns exemplos de listas: A seguir uma lista numerada

\begin{enumerate}
\item Pão
\item Queijo
\item Presunto
\end{enumerate}

A seguir uma lista não numerada

\begin{itemize}
\item Maria
\item João
\end{itemize}



\section{Tabelas}
Exemplo de tabela. A Tabela \ref{tab:comparacao} faz uma comparação de diversos parâmetros.

\begin{table}[h]
\begin{center}
\caption{Tabela de comparação.}
\label{tab:comparacao}
\begin{tabular}{|l|l|r|l|}
\hline
lattice & $d$ & $q$ & $T_{\rm mf}/T_c$ \\
\hline
square & 2 & 4 & 1.763 \\
\hline
triangular & 2 & 6 & 1.648 \\
\hline
diamond & 3 & 4 & 1.479 \\
\hline
simple cubic & 3 & 6 & 1.330 \\
\hline
bcc & 3 & 8 & 1.260 \\
\hline
fcc & 3 & 12 & 1.225 \\
\hline
\end{tabular}
\end{center}
\end{table}





\section{Figuras}

A Figura \ref{fig:Grafico} tem sua largura igual a 70\% da largura do texto.

\begin{figure}[h!]
\centering
\includegraphics[width=0.70\textwidth]{Figuras/Plot1.pdf} % a figura esta localizada na pasta Figuras
\caption{Gráfico colorido}
\label{fig:Grafico}
\end{figure}

Figuras e tabelas flutuam no texto. O \LaTeX{} tenta posicioná-las de forma a otimizar o máximo de espaço.




\section{Equações}
Veja com é fácil introduzir equações com linha própria.
%
\begin{equation}
\Delta =\sum_{i=1}^N w_i (x_i - \bar{x})^2 .
\end{equation}

É uma boa ideia numerar equações, Mas também podemos ter equações sem numeração.
%
\begin{equation*}
P(x) = \frac{x - a}{b - a} ,
\end{equation*}
%
and
%
\begin{equation*}
g = \frac{1}{2} \sqrt{2\pi} .
\end{equation*}

Usando o comando label é possível referenciar a equação depois.
%
\begin{equation}
\label{eq:ising}
E = -J \sum_{i=1}^N s_i s_{i+1} ,
\end{equation}

A equação~\eqref{eq:ising} expressa a energia no modelo Ising.

É possível definir novo comandos, veja no preâmbulo logo após os usepackages. Foram definidos os comandos Dx, Dy, imax e jmax:
%
\begin{gather}
\Dx = \frac{L}{\imax} \\
\Dy = \frac{H}{\jmax}
\end{gather}


Exemplos de outras equações:
%
\begin{equation}
I = \! \int_{-\infty}^\infty f(x)\,dx \label{eq:fine}.
\end{equation}

A equação \eqref{eq:fine} tem uma integral\footnote{Você pode inserir uma nota de rodapé facilmente.}.

pode-se alinhar equações:
%
\begin{align}
a + d & = b \\
c &= d +f ,
\end{align}
ou usar sub-equações:
\begin{subequations}
%
\begin{gather}
a  = b \\
c = d .
\end{gather}
\end{subequations}

Equações muito longas podem ser quebradas em mais de uma linha com o comando multline:
%
\begin{multline}
\label{Eq:Eq3}
\sum_{n=1}^{\infty} n = 1+2+3+4+5+6+7+8+9+10+11
\\
+12+13+14+15+16+17+18+19+20+\dots
\end{multline}


podemos escrever matrizes
%
\begin{align}
\textbf{T} &=
\begin{pmatrix}
T_{++} \hfill & T_{+-} \\
T_{-+} & T_{--} \hfill
\end{pmatrix} , \nonumber \\
& =
\begin{pmatrix}
e^{\beta (J + B)} \hfill & e^{-\beta J} \hfill \\
e^{-\beta J} \hfill & e^{\beta (J - B)} \hfill
\end{pmatrix}.
\end{align}

Abaixo, temos mais um exemplo de equação.
%
\newcommand{\rv}{\textbf{r}}
\begin{equation}
\sum_i \vec A \cdot \vec B = -P\!\int\! \rv \cdot
\hat{\mathbf{n}} \, \mathrm{d} A = P\!\int \! {\vec \nabla} \cdot \rv \, \mathrm{d} V.
\end{equation}



\section{Simbolos especiais}

\subsection{Letras gregas mais comuns}

Aqui temos algumas letras gregas:
$\alpha,
\beta, \gamma, \Gamma,
\delta,\Delta,
\epsilon, \zeta, \eta, \theta, \Theta, \kappa,
\lambda, \Lambda, \mu, \nu,
\xi, \Xi,
\pi, \Pi,$

Mais letras gregas
$
\rho,
\sigma,
\tau,
\phi, \Phi,
\chi,
\psi, \Psi,
\omega, \Omega$

\subsection{Símbolos especiais}

A função \emph{gather} permite incluir diversas equações de uma vez. A derivada é definida como
%
\begin{gather}
\frac{\mathrm{d}y}{\mathrm{d}x} = \lim_{\Delta x \to 0} \frac{\Delta y}
{\Delta x}
\\
f(x) \to y \quad \mbox{as} \quad x \to
x_{0}
\\
f(x) \mathop {\longrightarrow}
\limits_{x \to x_0} y
\end{gather}

Ordem de magnitude
%
\begin{gather}
\log_{10}f \simeq n
\\
f(x)\sim 10^{n}
\end{gather}
Igualdade de aproximação:
\begin{gather}
f(x)\simeq g(x)
\\
f(x) \propto x^3 .
\end{gather}


\end{document} 